% Post-checkpoint-14 auxiliary refinement; no representation-count conclusion.
\begin{proposition}[Proved: weighted compression of the exceptional graph]
\label{D:weighted-exceptional-chain-compression}
Retain the notation and hypotheses of
Proposition~\ref{D:exceptional-degree-linear-core}, including the
degree $n$ and constants $a,b,c=3b+a$.  For a component $F$ of the
exceptional support, let $m_F$ be its positive integral coefficient
in $f^*A$.  Mark a noncore vertex if its curve meets the core.
Then
\[
 \sum_{F\in\mathcal K}m_F(-F^2)\le nc,
 \qquad
 \sum_{Q\text{ marked}}m_Q\le nc.                    \tag{D.WC.1}
\]
Every connected noncore block meets the core, so there are at most
$nc$ such blocks.

Each noncore block is a path, or a tree with one trivalent vertex.
In the latter case its three arm lengths are either $(1,1,h)$,
with $h\ge1$, or $(1,2,h)$, with $h\in\{2,3,4\}$.  Let $P$ consist
of the marked vertices together with all trivalent vertices.
Then
\[
                 \#P\le2nc,
                 \qquad \sum_{Q\in P}m_Q\le31nc.      \tag{D.WC.2}
\]
Add a virtual endpoint of weight zero one edge beyond each unmarked
terminal leaf.  Decompose the graph into maximal segments whose
endpoints belong to $P$ or are these virtual endpoints.  The
weights on each segment form an integral arithmetic progression.
The sum of the lengths in edges of all nonconstant segments is
at most $93nc$.  There are at most $3nc$ constant segments, all
with two endpoints in $P$ and the same positive weight at both
ends.  On replacing each constant segment by one edge, the number
of remaining actual noncore vertices is at most $95nc$.

Thus the only unrestricted graph data in this estimate are the
lengths of at most $3nc$ constant-multiplicity rational
$(-2)$-chains.  The proposition does not assert a bound for those
lengths, an object atlas, or a representation-count consequence.
\end{proposition}

\begin{proof}
Use the nonnegative budget from (D.GC.4).  Write its dominating
and positive point-mapped contributions as $B_1$ and $B_2$;
thus $B_1+B_2\le nb$.  The adjunction inequalities used there
give
\[
 \sum_{F\in\mathcal K}m_F(-F^2)
 \le B_1+\sum_{F\text{ dominating}}m_F+3B_2
 \le3nb+na=nc.                                      \tag{D.WC.3}
\]
The pullback cycle is anti-nef on every component of its support:
projection gives $(f^*A)\cdot F<0$ for a dominating component and
$(f^*A)\cdot F=0$ for a point-mapped one.  Therefore
\[
 \Phi:=\sum_{F\in\mathcal K,\ Q\notin\mathcal K}
                         (F\cdot Q)m_Q
 \le\sum_{F\in\mathcal K}m_F(-F^2)\le nc.           \tag{D.WC.4}
\]
Every marked vertex occurs in this sum at least once, proving
the second assertion of (D.WC.1).

Let $N$ be the negative intersection matrix of a connected
noncore block, and let $T$ be its incidence matrix with the core.
Intersection with its point-mapped components gives
\[
                           N m=T^{\mathsf t}m_{\mathcal K}.
\]
If the block had no attachment, positive definiteness of $N$
would give $m=0$, contrary to the positivity of the coefficients
of $f^*A$.  Each block thus contains a marked vertex, and the
number of blocks is at most $\sum_{Q\text{ marked}}m_Q\le nc$.

We record the elementary graph argument, rather than using a
classification of surface singularities.  The matrix is $N=2I-P_0$
with $P_0$ symmetric, nonnegative and integral, and
$\rho(P_0)<2$.  A multiple edge of multiplicity at least two
already violates positive definiteness on its two vertices.
A cycle has adjacency spectral radius two, as does a four-arm
star.  Hence the graph is a tree of valency at most three.
If it contained two trivalent vertices, take two with no
trivalent vertex between them, their joining path, and the
first vertex on each of the four remaining incident edges.
On this principal subgraph, put weight two on the joining path
and weight one on those four leaves.  This is a positive
eigenvector of its adjacency matrix with eigenvalue two,
again impossible.  There is at most one trivalent vertex.

If the arm lengths at that vertex are $u\le v\le w$, elimination
of the three path matrices gives the positive Schur complement
condition
\[
                 \frac1{u+1}+\frac1{v+1}+\frac1{w+1}>1.
                                                               \tag{D.WC.5}
\]
It forces $u=1$.  Either $v=1$, or $v=2$ and $w\in\{2,3,4\}$.
This proves the asserted list by elementary matrix algebra.

There is at most one trivalent vertex per block.  Consequently
$\#P\le\#\{\text{marked vertices}\}+\#\{\text{blocks}\}\le2nc$.
It remains to bound the weight of an unmarked trivalent vertex.
Let its weight be $z$.  On each arm stop at the first marked
vertex, if there is one.  Otherwise stop at the virtual zero
vertex one edge beyond the terminal leaf.  Write $L_i$ for the
distance to the stop and $v_i$ for its weight.  At an unmarked
vertex the pullback intersection equation is
\[
                    2m_Q=\sum_{Q'\sim Q}m_{Q'}.
                                                               \tag{D.WC.6}
\]
Thus weights vary affinely up to the stops, and at the central
vertex (D.WC.6) reads
\[
             z\left(\sum_{i=1}^3\frac1{L_i}-1\right)
                       =\sum_{i=1}^3\frac{v_i}{L_i}.             \tag{D.WC.7}
\]
The coefficient on the left is positive by (D.WC.5), since a
stop can only shorten an arm.  For arms $(1,1,h)$, if neither
short arm contains a mark, then $L_1=L_2=2$, $v_1=v_2=0$, and
(D.WC.7) gives $z=v_3$.  Otherwise one of $L_1,L_2$ is one and
the coefficient is at least $1/2$, giving $z\le2\sum_i v_i$.
For arms $(1,2,h)$, $h\in\{2,3,4\}$, the coefficient is at least
\[
                         \frac12+\frac13+\frac15-1=\frac1{30}.
\]
Hence $z\le30\sum_i v_i$.  These first marked vertices are
distinct within a block and belong to no other block.  Adding
the already marked vertices proves
$\sum_{Q\in P}m_Q\le31\sum_{Q\text{ marked}}m_Q\le31nc$.

After adjoining the specified virtual endpoints, every actual
vertex outside $P$ has valency two.  Equation (D.WC.6) says that
successive weight differences on a maximal segment are equal
integers.  If a segment of length $\ell$ is nonconstant, that
integer is nonzero, so
\[
              \ell\le|m_{\rm right}-m_{\rm left}|
                         \le m_{\rm right}+m_{\rm left}.
\]
Each vertex of $P$ is an endpoint of at most three segments,
and virtual endpoints contribute zero weight.  Summing gives
\[
                \sum_{\text{nonconstant segments}}\ell
                    \le3\sum_{Q\in P}m_Q\le93nc.               \tag{D.WC.8}
\]
A constant segment cannot end at a virtual zero vertex, since
all actual weights are positive.  Its two endpoints lie in
$P$.  The number of constant segments is at most
$3\#P/2\le3nc$.  Once their internal vertices are removed, the
number of remaining actual vertices is at most
$\#P+\sum_{\text{nonconstant}}\ell\le95nc$, as claimed.
\end{proof}
